%% =============================================================================
%% INTRODUCTION
%% =============================================================================

\section{Introduction} \label{sec:intro}

The idea of talking to a computer to analyze data is not new. Natural
language interfaces to databases date to the
1970s~\citep{Woods:1973, Codd:1974}, and every generation of statistical software
has tried to lower the barrier between analytical intent and software
commands. What \emph{is} new is that the idea suddenly works well
enough to be everywhere. Since late 2022, millions of users have been
uploading datasets to ChatGPT and typing requests like ``compare last
quarter's sales by region''~\citep{Hu:2023}. The conversational
interface to data analysis has gone from research curiosity to
mass-market product in under two years.

The results are concerning. The model generates a \proglang{Python}
script and executes it on a remote virtual machine---the user's data
now resides on a third-party server they do not control. Even when
users copy generated code back to their own machine, the problems
persist: generated code may compute the wrong quantity without raising
an error~\citep{Ludwig+Mullainathan+Rambachan:2025}, and the same
prompt yields different code as models update---making results
irreproducible. Section~\ref{sec:defining} examines these failure
modes in detail.

The core idea---natural language as interface---is sound. The
architecture is not. These problems are not inherent to
conversational data analysis; they arise from a specific design
choice: delegating computation to generated code. This paper describes
an alternative approach we call \emph{chat-first data analytics}, built
on three principles:

\begin{enumerate}
  \item \textbf{Interpretation and computation are separated.}
    The language model selects methods and extracts parameters from natural
    language; pre-implemented, validated algorithms perform the actual
    computation.
  \item \textbf{Methods are pre-validated, not generated.}
    Correctness is established once, through systematic validation against
    reference implementations, rather than hoped for anew with each request.
  \item \textbf{Data stays local; only structure travels.}
    The language model receives column names, data types, and summary
    statistics---not the data itself. For complete isolation, the language
    model can run locally too.
\end{enumerate}

These principles address each concern directly. Correctness follows
from using validated implementations rather than generated code.
Privacy is protected by keeping data on the user's machine.
Reproducibility is achieved because the same tool invocation produces
identical results---the stochastic element of the language model
affects only interpretation, not computation.

The target users are researchers and analysts who know what analysis
they want but prefer not to write the code themselves: graduate
students running standard econometric models for the first time,
applied researchers iterating on specifications during exploratory
work, and data analysts in organizations where privacy regulations
prohibit cloud processing. For expert users fluent in \proglang{R}
or \proglang{Stata}, the system offers faster iteration on routine
tasks and a local privacy guarantee; for less experienced users, it
lowers the barrier to correctly specified analyses---though
statistical literacy remains necessary to interpret results
(Section~\ref{sec:defining}).

This paper presents \pkg{prompt2analytics} (\pkg{p2a}), an implementation of
chat-first data analytics. The software includes an econometrics and
data analytics library in pure \proglang{Rust}, exposed through the
Model Context Protocol (MCP) for integration with language models. The
implementation covers ordinary least squares (OLS) with robust standard errors,
panel data models with fixed effects (FE) and random effects (RE),
instrumental variables (IV), difference-in-differences (DiD),
discrete choice, time series, and over 200 additional methods---all
validated against established \proglang{R} packages.

\pkg{prompt2analytics} is available as open-source software under the MIT
license at \url{https://github.com/umatter/prompt2analytics}.
Supplementary materials---including all 96 evaluation prompts,
benchmark scripts for both \proglang{R} and \proglang{Rust}, and validation
test suites---are located in the repository's \code{paper/} directory.
Appendix~\ref{app:reproducibility} provides reproduction instructions.

The remainder of the paper is organized as follows.
Section~\ref{sec:defining} motivates chat-first analytics by
examining the failure modes of current approaches and situating the
three principles within existing work.
Section~\ref{sec:tools} describes the software implementation.
Section~\ref{sec:examples} demonstrates the system through worked
examples including realistic failure cases.
Section~\ref{sec:evaluation} evaluates tool selection accuracy and
numerical correctness.
Section~\ref{sec:deployment} presents performance benchmarks and local
deployment configurations, and Section~\ref{sec:discussion} discusses
limitations and future directions.
